% =====================================================================
% Draft-only task list. Rendered as a red box on the title page while
% \showtaskstrue; hidden by \showtasksfalse. Not part of the manuscript
% and excluded from the word count.
% =====================================================================
\label{sec:tasks-open}

\begin{center}
\fbox{\parbox{0.94\textwidth}{%
\textbf{\textcolor{red}{OPEN ITEMS BEFORE SUBMISSION}}\\[.4em]
\footnotesize

\textbf{A. Cross-references to companion papers}\\
Three placeholder bib keys are used throughout:
\texttt{oded2026intervene} (the AAAI methodological paper),
\texttt{oded2027kineret-jbi} (the JBI adherence paper), and
\texttt{mediator-github}. Update these to the accepted / preprint form
once the two companion drafts are placed. The \emph{Background} section
also references the JBI paper's Table~1 for cohort characteristics; if
JBI is not yet publicly citeable, reproduce the table as an appendix
here.\\[.2em]

\textbf{B. Empirical additions}\\
The following extensions are being staged in the Kineret-Prediction
repository. Items marked \emph{[inference-only]} reuse the existing
per-arm checkpoints and do not require retraining.
\begin{itemize}\setlength{\itemsep}{.05em}
  \item \emph{[inference-only]} Sensitivity of the ladder ordering to
        $K_{\mathrm{eval}} \in \{2, 3, 5, 6\}$. Because training augments
        over $K \in 1\ldots 13$, each trained checkpoint can be re-scored
        at any $K_{\mathrm{eval}}$ by cutting the same test admissions at
        a different point; no retraining is required.
  \item \emph{[inference-only]} Reliability curves and per-outcome ECE
        (with optional temperature scaling on validation), as reported
        for INTERVenE-Enc on MIMIC-IV in \citep{oded2026intervene}.
        Predictions live in \texttt{results/full\_bench\_results/runs/<arm>/}.
  \item \emph{[inference-only]} Age-band and sex subgroup AUROC/AUPRC
        (item C below); paired-bootstrap $\Delta$ and BH-FDR across
        outcomes (item D). Both from the existing predictions.
  \item \emph{[requires re-run]} Temporal held-out: retrain the ladder
        on pre-guideline admissions and evaluate on post-guideline (or
        vice versa), to test whether prediction quality tracks the
        practice shift documented in the JBI companion paper. Supported
        natively by the pipeline via \texttt{date\_range\_start} /
        \texttt{date\_range\_end}.
  \item \emph{[requires re-run]} Per-hospital leave-one-out for the KB
        arm (six re-runs), to quantify site generalization. Requires the
        hospital column to be plumbed through cohort assembly; check
        \texttt{Kineret-Prediction} PR before running.
  \item \emph{[inference-only]} Sparse-transcoder attribution
        reproduction on the KB+QA arm on this cohort. A partial run
        exists under \texttt{results/transformers-results-old-with-xai/}
        for pre-ladder ss-STraTS variants; the KB+QA reproduction reuses
        the pipeline documented in \citep{oded2026intervene}.
  \item \emph{Optional.} Multi-seed variance for the KB and KB+QA arms.
        The AAAI companion paper explicitly declined multi-seed at this
        cohort scale, on cost grounds, and reported patient-level
        bootstrap intervals only; we adopt the same stance and mention
        it as a limitation.
\end{itemize}
\vspace{.1em}

\textbf{C. Fairness / subgroup analyses}\\
AIIM reviewers routinely ask for at least age-band and sex subgroup
AUROC/AUPRC. The cohort is $54.3\%$ male and skewed older; a
per-subgroup table for the KB arm is one bootstrap run away.\\[.2em]

\textbf{D. Statistical practice}\\
\begin{itemize}\setlength{\itemsep}{.05em}
  \item Decide whether to report $P$ values for pairwise arm comparisons
        (currently only overlap-of-CIs is used). If so, the paired
        bootstrap $\Delta$ distribution is already computed for
        Figure~\ref{fig:qa-delta}; extend to all pairs.
  \item Multiple-testing correction across the eight per-outcome
        comparisons in Table~\ref{tab:per-outcome} (BH-FDR at $q<0.05$
        matches the JBI paper's convention).
\end{itemize}
\vspace{.1em}

\textbf{E. Author contributions and funding}\\
Funding statement outstanding (aligned with JBI/JAMIA companion
papers); ORCID identifiers to be added for all authors.\\[.2em]

\textbf{F. Figure budget}\\
AIIM allows up to $\sim 10$ figures without special permission. The
current draft uses $5$. Candidates for expansion: (i) event-reconciliation
diagram of the raw-vs-Mediator alignment pipeline (methods section);
(ii) sparse-transcoder attribution panel for the KB+QA arm;
(iii) per-outcome AUROC alongside AUPRC (data exists as
\texttt{fig4\_per\_target\_auroc.pdf} but is not included in the current
draft). Add if space permits after the empirical additions in item B.\\[.2em]

\textbf{G. Language and prose}\\
The draft borrows phrasing from the concurrent JBI/JAMIA drafts by
design (shared methods sections). Before submission, run a self-plagiarism
check against those two documents; substantive text overlap should
either be paraphrased or attributed as a shared methods description.
The abstract is currently unstructured. AIIM does not require a
structured abstract but common editorial policy is to prefer one; consider
converting to Objective / Methods / Results / Conclusion form.\\[.2em]

\textbf{H. Compilation}\\
Not yet built --- no TeX distribution here. Before submission, compile
end-to-end with \texttt{pdflatex} + \texttt{bibtex} and reconcile any
line-breaking / hyphenation issues in the tables. Confirm figure
placement floats do not push results captions past their referring
paragraph.
}}
\end{center}

\clearpage
